\documentclass[10pt,a4paper]{amsart}
\usepackage[margin=19mm]{geometry}
\usepackage{amsmath,amssymb,mathtools}
\usepackage[scr=rsfs,cal=euler]{mathalfa}
\usepackage{xfrac}
\usepackage[unicode,hidelinks,bookmarks=false]{hyperref}
\newcommand{\ie}{i.e.\ }
\newcommand{\cf}{cf.\ }
\newcommand{\resp}{respectively\ }
\newcommand{\Subsection}{\subsection}
\providecommand{\nobreakdash}{\nobreak\hbox{-}}
\setlength{\emergencystretch}{2em}
\multlinegap0pt
\usepackage{bm}
\usepackage{bbm}
\usepackage{dsfont}
\usepackage{xspace}
\usepackage{cancel}
\usepackage{mathtools}
\usepackage{amsthm}
\makeatletter
\@ifundefined{theorem}{\newtheorem{theorem}{Theorem}}{}
\makeatother
\newcommand{\R}{\mathbb{R}}
\newcommand{\bs}{\setminus}
\title[A note on irreducibility for topical maps]{A note on irreducibility for topical maps}
\author{Brian Lins}
\date{Source: arXiv:2510.12522v1 (CC BY 4.0)}
\begin{document}
\maketitle

\noindent\emph{Source and licence.} A note on irreducibility for topical maps. Version arXiv:2510.12522v1 (CC BY 4.0), distributed under the Creative Commons Attribution 4.0 International licence, as recorded in the arXiv licence metadata for this exact version and shown on the abstract page https://arxiv.org/abs/2510.12522v1. Full rights notice: licenses/arxiv-2510.12522-CC-BY-4.0.txt.\par\smallskip

\noindent\textit{Editorial scope.} Counted unit: the Theorem whose source label is \texttt{thm:connect1} in the article (source0.tex lines 245--254, source label \texttt{thm:connect1}), delivered together with the article's own complete original proof of it (source0.tex lines 256--290, one \texttt{proof} environment of 35 lines). No mathematics is written, repaired or re-derived: statement and proof are the article's own text line for line. Required context retained uncounted: eq:face (lines 161--163). Every definition, display and label of those lines is retained unchanged. Omitted from this package and not redistributed: 1--160, 164--244, 255--255, 291--811. Nothing inside a retained excerpt is omitted or altered: every retained body is delivered exactly as the article's lines are written, so no omission receipt of any kind accompanies this package. Citations: the retained excerpts carry no citation command at all, so no bibliography entry of the article is transcribed and no reference is redistributed. Reference adaptation: none was needed and none was made; every reference inside the delivered unit resolves inside it, so no unresolved reference or citation is produced. Printed slips kept literal: no printed word, symbol, hypothesis or number of the counted statement or of its proof was altered. Layout and class: the article's own class is \texttt{amsart} with packages including amsmath, amsfonts, amsthm, amssymb, enumitem, tikz, soul, hyperref; the wrapper is the lane's pinned \texttt{amsart} editorial wrapper at 10pt on A4, which restates the theorem-like environments the retained text needs and adds the article's own macro definitions that the retained text needs (\texttt{\textbackslash R}, \texttt{\textbackslash bs}), with extra wrapper packages (bm, bbm, dsfont, xspace, cancel, mathtools). The delivered numbering is the unit's own. Assets: no retained span contains a figure environment, a graphics-inclusion command, a PDF-page inclusion, a table or a TikZ picture; no image file is copied into this package, the package declares no asset dependency and no image file is redistributed with it. Licence: version arXiv:2510.12522v1 of this article is distributed under the Creative Commons Attribution 4.0 International licence, as recorded in the arXiv licence metadata for this exact version and shown on the abstract page https://arxiv.org/abs/2510.12522v1. A full rights notice is delivered at licenses/arxiv-2510.12522-CC-BY-4.0.txt. Characters: every retained excerpt is pure ASCII, so the delivered engine, XeLaTeX with its default Latin Modern text font, reports no missing character.
\begin{equation} \label{eq:face}
F_J := \{ x \in \R^n_{\ge 0} : x_i = 0 \text{ for all } i \notin J \}
\end{equation}

\begin{theorem} \label{thm:connect1}
Let $f:\R^n_{\ge 0} \rightarrow \R^n_{\ge 0}$ be an m-topical map.  
\begin{enumerate}
\item \label{item:face=>indecomp} If $f$ is facially or graphically irreducible, then $f$ is indecomposable. 
%\item \label{item:graph=>indecomp} If $f$ is graphically irreducible, then $f$ is indecomposable. 
\item \label{item:mconvex1} If $f$ is m-convex and indecomposable, then $f$ is graphically irreducible.  
%\item \label{item:mconvex2} If $f$ is m-convex and facially irreducible, then $f$ is graphically irreducible.  
\item \label{item:subadditive} If $f$ is subadditive and indecomposable, then it is facially irreducible. %and graphically irreducible.
\end{enumerate}
\end{theorem}

\begin{proof}
Proof of \eqref{item:face=>indecomp}.  Let $J$ be any nonempty proper subset of $[n]$. If $f$ is facially irreducible, there must be an $i \in [n] \bs J$ such that $f(e_J)_i > 0$. Then 
$$\lim_{t \rightarrow \infty} f(\mathbf{1} + t e_J)_i \ge \lim_{t \rightarrow \infty} t f(e_J)_i = \infty.$$
So we conclude that $f$ is indecomposable.  

%~\\
%Proof of \eqref{item:graph=>indecomp}. 
If $G(f)$ is strongly connected, then for any nonempty proper subset $J \subset [n]$, there must be an edge from some $i \notin J$ to some $j \in J$.  So 
$$\lim_{t \rightarrow \infty} f(\mathbf{1} + t e_J)_i \ge \lim_{t \rightarrow \infty} f(\mathbf{1} + t e_{\{j\}})_i = \infty.$$
Therefore $f$ is indecomposable. 

~\\
Proof of \eqref{item:mconvex1}. 
Suppose that $f$ is m-convex and indecomposable, but $G(f)$ is not strongly connected.  Let $I$ be a final class of $G(f)$ and let $J := [n] \bs I$.  Since $I$ is final, there are no edges from $i \in I$ to any $j \in J$.  Therefore 
$$\lim_{t \rightarrow \infty} f(\mathbf{1} + t e_{\{j\}})_i < \infty$$
for each $i \in I$ and $j \in J$.  
Observe that the entrywise geometric mean of the vectors $\{ \mathbf{1} + t e_{\{j\}} : j \in J\}$ is $\mathbf{1} + t^{1/|J|} e_J$. Since $f$ is m-convex, it follows that 
$$f(\mathbf{1} + t^{1/|J|} e_J)_i  \le \left( \prod_{j \in J} f(\mathbf{1} + t e_{\{j\}})_i \right)^{1/|J|}$$
for every $i \in I$. Therefore 
$$\lim_{t \rightarrow \infty} f(\mathbf{1} + t e_J)_i =\lim_{t \rightarrow \infty} f(\mathbf{1} + t^{1/|J|} e_J)_i  \le \left( \prod_{j \in J} \lim_{t \rightarrow \infty} f(\mathbf{1} + t e_{\{j\}})_i \right)^{1/|J|} < \infty.$$
Since this holds for all $i \notin J$, we have contradicted the assumption that $f$ is indecomposable, so we conclude that $G(f)$ must be strongly connected and therefore $f$ is graphically irreducible.  

%~\\
%Proof of \eqref{item:mconvex2}. 
%If $f$ is m-convex and facially irreducible, then it is also indecomposable by \eqref{item:face=>indecomp}.  So it is graphically irreducible by \eqref{item:mconvex1}. 

~\\
Proof of \eqref{item:subadditive}.
Suppose that $f$ is subadditive and indecomposable.  Consider any nonempty proper subset $J \subset [n]$. 
Since $f$ is indecomposable, there must be at least one $i \in [n] \bs J$ such that 
$$\lim_{t \rightarrow \infty} f(\mathbf{1} + te_J)_i  = \infty.$$
Then by subadditivity, 
$$\lim_{t \rightarrow \infty} f(\mathbf{1})_i + tf(e_J)_i  = \infty.$$
So we conclude that $f(e_J)_i > 0$, and therefore the face $F_J$ defined by \eqref{eq:face} cannot be invariant.  This proves that $f$ is facially irreducible. % Since subadditive m-topical functions are also m-convex, it follows from \eqref{item:mconvex1} that $f$ is also graphically irreducible.
\end{proof}

\end{document}
